\subsubsection{Etapa 2.3: re-reclutamiento y sustitución tras un fallo}
\label{subsubsec:sp2-recovery}

Cuando falla un miembro, la coalición puede perder el certificado de soporte, fuerza o torque aun con la carga detenida. E6 publica ese déficit y permite que las reservas revisen su participación antes de reanudar la misión. La campaña considera una carga y un fallo; el reemplazo viaja a velocidad acotada hacia un contacto prefijado en una trayectoria libre.

\paragraph{Fallo y reserva}

El reemplazo solo se moviliza después de que las revisiones locales reduzcan $D_6(\boldsymbol x)$ a cero, y la guardia vuelve a comprobar el certificado antes del arribo (Figura~\ref{fig:sp2-compact-pipeline}).

% [figure omitido aquí: figura didáctica ya en el cuerpo; ver cuerpo, Sección 6]

\paragraph{Problema de optimización de referencia}

Sea $\mathcal R_k$ la reserva, $n_R=|\mathcal R_k|$, $A_6$ los recursos, $\boldsymbol d_k^6$ el déficit y $\boldsymbol c_i^6$ la contribución normalizada. Sus tres componentes son soporte, fuerza y torque. $x_i$ recluta y $\kappa_i^6$ agrega desplazamiento y cambio adimensionales.

El oráculo central resuelve
\begin{equation}
\label{eq:sp6-oracle}
  \min_{\boldsymbol x\in\{0, 1\}^{n_R}}
  \sum_{i\in\mathcal R_k}\kappa_i^6x_i
  \quad\text{sujeto a}\quad
  \sum_{i\in\mathcal R_k}c_{ia}^6x_i\geq d_{ka}^6,
  \quad a=1,\ldots,A_6.
\end{equation}
La enumeración de los $2^{n_R}$ perfiles certifica el coste mínimo postfallo con información global y fija el techo para los métodos locales.

\paragraph{Métodos de comparación}

ALLIANCE aporta el antecedente de activación alternativa \parencite{parker1998alliance}; las subastas dinámicas, ACBBA y CBBA-PR motivan la revisión asíncrona tras eventos \parencite{nanjanath2006dynamicAuctions,johnson2011acbba,buckman2019cbbaPr}. Como esas fuentes no incorporan el certificado de \emph{wrench} usado aquí, la Tabla~\ref{tab:sp2-compact-e6-methods} distingue el mecanismo de reasignación de la guardia física.

% [table omitido aquí: tabla de métodos ya en el cuerpo; ver cuerpo, Sección 6]

\paragraph{Juego de reparación}

Defínanse el déficit residual y el coste de movilización
\begin{equation}
\label{eq:sp6-deficit-potential}
\begin{aligned}
 D_6(\boldsymbol x)
 &=\sum_{a=1}^{A_6}w_a
 \left[d_{ka}^6-\sum_{i\in\mathcal R_k}c_{ia}^6x_i\right]_+,\\
 K_6(\boldsymbol x)&=\sum_i\kappa_i^6x_i,
 &\Phi_6(\boldsymbol x)&=-\lambda_6D_6(\boldsymbol x)-K_6(\boldsymbol x).
\end{aligned}
\end{equation}
En~\eqref{eq:sp6-deficit-potential}, $w_a>0$ y $\lambda_6>0$, y todo término es adimensional. El robot usa una utilidad de contribución marginal,
\begin{equation}
\label{eq:sp6-payoff}
 U_i(x_i,\boldsymbol x_{-i})=
 \Phi_6(x_i,\boldsymbol x_{-i})-\Phi_6(0,\boldsymbol x_{-i}).
\end{equation}
La carga publica déficit y cada reserva usa $\boldsymbol c_i^6,\kappa_i^6$. Calibrar $\lambda_6$ exactamente usa la instancia completa fuera de línea.

\begin{teorema}[potencial exacto y terminación]
\label{thm:sp6-exact-potential}
El juego definido por~\eqref{eq:sp6-payoff} es un juego de potencial exacto con potencial $\Phi_6$. Existe al menos un Nash puro y todo camino maximal de mejoras estrictas termina en un Nash después de, como máximo, $2^{n_R}-1$ cambios aceptados.
\end{teorema}

$\boldsymbol x$ es Nash si ningún robot libre satisface
$\lambda_6[D_6(\boldsymbol x)-D_6(\boldsymbol x+\boldsymbol e_i)]>\kappa_i^6$
y ningún miembro ahorra al salir más que el déficit inducido. Estas dos condiciones caracterizan todos los equilibrios; la derivación aparece en el Anexo~\ref{app:sp6-proofs}.

Para excluir equilibrios deficitarios, sea
\begin{equation}
\label{eq:sp6-delta-min}
 \delta_{min}=min_{\boldsymbol x:D_6(\boldsymbol x)>0}
 \max_{j:x_j=0}\left[D_6(\boldsymbol x)-D_6(\boldsymbol x+\boldsymbol e_j)\right],
 \qquad
 \kappa_{max}^6=\max_i\kappa_i^6.
\end{equation}

\begin{teorema}[caracterización exacta de los Nash]
\label{thm:sp6-feasible-nash}
Supónganse capacidades no negativas, $w_a>0$, costes $\kappa_i^6>0$ para todo $i$ y un déficit postfallo no trivial $D_6(\boldsymbol0)>0$. Si la reserva completa cubre $\boldsymbol d_k^6$, entonces $\delta_{\min}>0$. Si $\lambda_6>\kappa_{max}^6/\delta_{\min}$, los Nash puros coinciden exactamente con las coaliciones factibles mínimas por inclusión:
\[
 \operatorname{NE}(\mathcal G_6)=
 \{\boldsymbol x:D_6(\boldsymbol x)=0,
   \boldsymbol x\text{ es mínima por inclusión}\}.
\]
\end{teorema}

La demostración considera las dos desviaciones unilaterales: ingresar reduce un déficit y salir de una coalición mínima crea al menos $\delta_{\min}$. Con el umbral indicado, la penalización domina el ahorro de salida y un ingreso redundante aumenta el coste (Anexo~\ref{app:sp6-feasible-nash-proof}).

Como los costes son positivos, toda reparación factible de coste mínimo es mínima por inclusión y, por el Teorema~\ref{thm:sp6-feasible-nash}, es Nash. De esta inclusión se sigue
\begin{equation}
 \operatorname{PoS}_6=1.
 \label{eq:sp6-pos}
\end{equation}
El conjunto de Nash contiene un óptimo social; la trayectoria de mejoras puede terminar en cualquier reparación mínima por inclusión.

\begin{proposicion}[ineficiencia no acotada del peor Nash]
\label{prop:sp6-unbounded-poa}
No existe una cota uniforme, independiente de la instancia y de $\lambda_6$, para el cociente entre el coste del peor Nash factible y el oráculo~\eqref{eq:sp6-oracle}, incluso con un solo recurso y tres robots.
\end{proposicion}

La construcción del Anexo~\ref{app:sp6-proofs} usa un soporte de coste $M$ y dos complementarios de coste total dos. Ambos son Nash y el cociente $M/2$ diverge, de donde $\operatorname{PoA}_6=+\infty$ aunque exista un equilibrio óptimo.

\begin{contribucion}{equilibrio de recuperación sujeto a certificado}
\estadoaporte{demostrado para el juego binario con déficit y costes congelados.}
La utilidad enlaza el déficit físico con la decisión de cada reserva. El umbral del Teorema~\ref{thm:sp6-feasible-nash} excluye equilibrios deficitarios y los costes positivos excluyen miembros redundantes. Estas condiciones no controlan el coste del peor equilibrio, que puede crecer sin cota frente al oráculo.
\end{contribucion}

\paragraph{Movimiento del reemplazo}

Una decisión guardada no reanuda instantáneamente el transporte. Para el reemplazo $i$, sean $\ell_i$ la distancia al contacto, $v_i^{\min}>0$ una velocidad garantizada y $\tau_s$ el asentamiento. Sea $\bar\tau_d$ una cota desde el fallo hasta su detección y $\bar\tau_a$ una cota para cada ventana agregada posterior: desde la detección o el último cambio aceptado hasta el siguiente cambio aceptado, o hasta la certificación terminal; la ventana incluye todas las activaciones que no cambian el perfil. Esta hipótesis cuantitativa es más fuerte que la sola justicia.

\begin{corolario}[cota temporal de reparación]
\label{cor:sp6-time-bound}
Bajo los supuestos de los Teoremas~\ref{thm:sp6-exact-potential}--\ref{thm:sp6-feasible-nash}, una política que prolonga un camino maximal y acota por $\bar\tau_a$ cada ventana agregada entre cambios aceptados o certificación, con trayectorias de reemplazo libres y velocidades inferiores positivas,
\begin{equation}
\label{eq:sp6-time-bound}
 T_{\mathrm{rec}}
 \leq \bar\tau_d+2^{n_R}\bar\tau_a
 +\max_{i:x_i=1}\frac{\ell_i}{v_i^{\min}}+\tau_s.
\end{equation}
Si el lado derecho no supera el tiempo seguro remanente $T_{\mathrm{safe}}$, el certificado puede restaurarse antes del plazo bajo este modelo de arribo.
\end{corolario}

El término $\ell_i/v_i^{\min}$ aproxima el viaje por un segmento libre y omite el giro y la evitación; $\tau_s$ representa el asentamiento por separado. La campaña evalúa cierre, restauración del certificado y puntualidad hasta el contacto.

\paragraph{Diseño experimental}

% sp6_numbers.tex ya se cargó en el cuerpo (sections/v2/sp2-compact.tex).
La Tabla~\ref{tab:sp6-design} resume \SPSixWorlds{} mundos pareados y \SPSixRuns{} ejecuciones. De ellos, \SPSixRecoverableWorlds{} disponen de reserva suficiente y \SPSixIrrecoverableWorlds{} carecen de un recurso; el endpoint principal es la restauración del certificado.

\begin{table}[!htbp]
  \centering
  \caption{Diseño experimental de E6-C.}
  \label{tab:sp6-design}
  \viutablefont
  \begin{tabularx}{0.94\textwidth}{>{\raggedright\arraybackslash}p{2.4cm} Y}
    \toprule
    Elemento & Especificación \\
    \midrule
    Factores & Cuatro escenarios, $n_R\in\{4,6,8\}$, semillas 8600--8639 y cinco métodos sobre cada mundo. \\
    Estimandos & Restauración, redundancia, coste frente al oráculo, puntualidad y mensajes. \\
    Inferencia & McNemar para proporciones, Wilcoxon para diferencias pareadas, IC \mbox{bootstrap} y corrección de Holm. \\
    Controles & Enumeración de perfiles, correspondencia de utilidad y potencial, condición de Nash, ascenso estricto, imposibilidad y cota temporal. \\
    \bottomrule
  \end{tabularx}
  \viusource{elaboración propia a partir del protocolo experimental de E6}
\end{table}

\paragraph{Resultados}

En los 480 mundos, todo método que intenta reparar alcanza una tasa de certificado de $0.75$, pues una cuarta parte de las instancias es imposible por construcción (Tabla~\ref{tab:sp6-results}).

\begin{table}[!htbp]
  \centering
  \caption{Resultados agregados de E6-C sobre 480 mundos pareados por método.}
  \label{tab:sp6-results}
  \scriptsize
  \resizebox{\textwidth}{!}{\begin{tabular}{lrrrr}
\toprule
Método & Cert. restaurado & Éxito temporal & Redundantes & Mensajes \\
\midrule
Juego potencial guardado & 0.750 & 0.544 & 0.000 & 24.0 \\
Subasta marginal & 0.750 & 0.596 & 0.175 & 86.4 \\
Greedy por distancia & 0.750 & 0.596 & 0.604 & 9.8 \\
Sin reparación & 0.000 & 0.000 & 0.000 & 0.0 \\
Oráculo exacto & 0.750 & 0.583 & 0.000 & 32.7 \\
\bottomrule
\end{tabular}
}
  \viuownsource
\end{table}

Condicionado a factibilidad física, el juego restaura el certificado en \SPSixConditionalRestore{} de los mundos y no reabre ninguno de reserva insuficiente. Frente a no reparar, el efecto no condicionado es \SPSixHOneEffect{}, y es determinista: en tres de los cuatro escenarios la reparación restaura el certificado en todas las semillas y no repararlo no lo restaura en ninguna, mientras que en \emph{infeasible reserve} ninguna de las dos lo consigue. El $0{,}750$ es, por tanto, la fracción de escenarios del banco en que la reparación actúa, idéntica en las 40~semillas; no procede publicar intervalo ni valor $p$ para un contraste sin variabilidad muestral. En la enumeración, todos los Nash fueron factibles y mínimos por inclusión, como predicen los teoremas.

El juego eliminó miembros unilateralmente redundantes, a diferencia del voraz, con efecto medio \SPSixHTwoEffect{} y $p_{\mathrm{Holm}}=\SPSixHTwoPHolm$. La reparación mínima por inclusión presentó una brecha media de movilización \SPSixPotentialGap{} frente al oráculo y contabilizó \SPSixPotentialMessages{} mensajes por mundo, entre el voraz y la subasta marginal. La diferencia absoluta de coste frente al oráculo fue \SPSixHThreeEffect{} ($p_{\mathrm{Holm}}=\SPSixHThreePHolm$).

En los escenarios recuperables sin plazo estrecho, todos los métodos activos restauraron el certificado y llegaron a tiempo. Al estrechar el plazo, el éxito temporal del juego cayó a \SPSixTightPotentialSuccess{} frente a \SPSixTightGreedySuccess{} del voraz; la distancia de la reserva domina una decisión estratégicamente factible cuando el tiempo físico resulta insuficiente.


No se definió una función que combinara coste, puntualidad y comunicación; el juego priorizó restauración sin redundancia y pagó una brecha de movilización frente al oráculo.

La validez depende sobre todo del certificado aditivo y de la calibración exponencial de $\delta_{\min}$ fuera de línea. Un QP geométrico y una trayectoria con obstáculos podrían cambiar tanto la coalición elegida como la puntualidad; los fallos múltiples no se ensayaron.

\paragraph{Síntesis}

Con reserva suficiente y una penalización por encima del umbral, los Nash puros coinciden con reparaciones factibles mínimas por inclusión y la mejor respuesta termina en un número finito de cambios. Las 480 instancias respetaron esa caracterización. El plazo estrecho reveló la limitación operativa: restaurar el certificado produjo éxito temporal \SPSixTightPotentialSuccess{}, por debajo de \SPSixTightGreedySuccess{} del voraz.

$\operatorname{PoS}_6=1$ garantiza que existe un Nash óptimo, pero $\operatorname{PoA}_6=+\infty$ permite equilibrios arbitrariamente caros. E7 toma las coaliciones restauradas como agentes compuestos y evalúa sus conflictos de ruta.
